\chapter{Analytic coefficients and collision monodromy}
\label{cbc:chap:analytic}

Formal translation is finite on each mode polynomial. Numerical
translation acts on an infinite series of modes and can destroy its
convergence. The first question is therefore quantitative: how much
jet radius does moving a point consume? For the Heisenberg
coefficients, a single series of bosons determines the exact loss.
The same seminorms then determine where the full line gauge exists.

A collision also changes the attachment of a coefficient to the
regular module. We retain this attachment in finite models of
punctured lines. Their higher products distinguish a marked
localization from its primitive quotient. Moving two punctures
together leads to an inverse telescope, while differentiation in
all spatial variables gives a different Hom complex. These
calculations provide explicit coefficient categories on which
collision transport can be studied.

We first work over $\mathbb C$, with $\kappa\ne0$. Differentials
have cohomological degree one, odd derivatives act on the left,
and exterior forms stand on the left of coefficient expressions.
The algebraic sections below work over a field $k$ of characteristic
zero. Every completion and every change of coefficient category
will be specified.

\section{Convergent jets and numerical translation}
\label{cbc:ac:completion}

For a finite set $I$, let
\[
 Y_I=\mathbb C[B_{i,n}:i\in I,\ n\geq0]
          \otimes\Lambda(C_{i,n}:i\in I,\ n\geq0),
 \qquad |B_{i,n}|=0,\quad |C_{i,n}|=1.
\]
Fix an order for the odd generators. For a monomial $m$, let
$d_i(m)$ be its number of factors at site $i$, counted with
multiplicity. Give $B_{i,n}$ weight $n$ and $C_{i,n}$ weight
$n+1$, and let $w_i(m)$ be the total weight at that site.
Field degree and jet weight will be controlled independently.

\begin{definition}\label{cbc:ac:seminorms}
For positive real numbers $A_i,r_i$, set
\[
 \left\|\sum_m f_m m\right\|_{\mathbf A,\mathbf r}
      =\sum_m |f_m|\prod_{i\in I}
                       A_i^{d_i(m)}r_i^{w_i(m)}.
\]
Given positive radii $\mathbf R=(R_i)_{i\in I}$, define
$\mathcal Y_I(\mathbf R)$ as the completion of $Y_I$ for
all these seminorms with $0<r_i<R_i$ and $A_i>0$.
Its cochain complex uses the homogeneous cohomological-degree
subspaces.
\end{definition}

Signed multiplication is submultiplicative for each seminorm.
A repeated odd generator gives zero. Taking rational amplitudes
cofinal toward infinity and rational radii cofinal below $R_i$
gives a countable defining family. Thus $\mathcal Y_I(\mathbf R)$
is a Fr\'echet algebra. It consists exactly of the coefficient
series for which every displayed absolute sum converges.
Polynomial truncations converge in this topology.

The subalgebra containing only $J_i=B_{i,0}$ is
$\mathcal O(\mathbb C^I_J)$, the entire functions of the
charges. Indeed its seminorms are the absolute power-series
sums at all positive amplitudes. The projection and inclusion are
\[
 \pi:\mathcal Y_I(\mathbf R)\longrightarrow
                           \mathcal O(\mathbb C^I_J)
 \quad\text{and}\quad
 \iota:\mathcal O(\mathbb C^I_J)\longrightarrow
                           \mathcal Y_I(\mathbf R).
\]

For an open position domain $U$, a holomorphic coefficient
series $\sum_m f_m(z)m$ is admitted when
\[
 \sum_m\sup_{z\in K}|f_m(z)|
             \prod_i A_i^{d_i(m)}r_i^{w_i(m)}<\infty
\]
for every compact $K\subset U$ and every defining seminorm.
Use these series as coefficients of holomorphic differential
forms, denoted by $\Omega^\bullet(U;\mathcal Y_I(\mathbf R))$.
Position derivatives are continuous by Cauchy estimates on
slightly larger compact subsets. Entire charge coefficients
use the same convention.

Recall the derivations of \eqref{cbc:cf:operators}:
\[
 QB_n=\kappa nC_{n-1},\quad QC_n=0,\qquad
 TB_n=nB_{n-1},\quad TC_n=nC_{n-1},\qquad
 LC_n=\kappa^{-1}B_n,\quad LB_n=0.
\]
An index $-1$ gives zero. At several sites take their sums.
If $N$ counts the ghosts and the positive-index bosons,
define
\begin{equation}\label{cbc:ac:h}
 h=N^{-1}\sum_{i,n\geq0}
        \frac{B_{i,n+1}}{\kappa(n+1)}\partial_{C_{i,n}}
       \quad(N>0),\qquad h=0\quad(N=0).
\end{equation}

\begin{proposition}\label{cbc:ac:continuous-contraction}
The operators $Q,T,L,h$ are continuous on
$\mathcal Y_I(\mathbf R)$. Their polynomial identities extend
to this completion. In particular,
\[
 Qh+hQ=1-\iota\pi,\qquad
 h^2=h\iota=\pi h=0,\qquad
 \|hF\|_{\mathbf A,\mathbf r}
              \leq|\kappa|^{-1}\|F\|_{\mathbf A,\mathbf r}.
\]
This is an $\mathcal O(\mathbb C^I_J)$-linear strong
deformation retract onto the entire charges in degree zero.
\end{proposition}
\begin{proof}
We give the one-site estimates. For $r<r'<R$, put
$c(r/r')=\sup_{w\geq0}w(r/r')^w$, which is finite.
The operator $Q$ preserves field degree and jet weight.
On a monomial of weight $w$, the sum of its absolute
coefficients is at most $|\kappa|w$. The operator $T$
lowers weight by one, with total coefficient at most $w$.
Consequently
\[
 \|QF\|_{A,r}\leq|\kappa|c(r/r')\|F\|_{A,r'},\qquad
 \|TF\|_{A,r}\leq r^{-1}c(r/r')\|F\|_{A,r'}.
\]

For $A<A'$, the operator $L$ has the bound
\[
 \|LF\|_{A,r}\leq
   \frac{\sup_{d\geq0}d(A/A')^d}{|\kappa|r}
                            \|F\|_{A',r}.
\]
These estimates prove continuity for the full seminorm family.
At finitely many sites the estimates add.

Every replacement in \eqref{cbc:ac:h} preserves field degree
and jet weight. Its absolute coefficient is
$1/(|\kappa|(n+1)N)$. The number of ghosts is at most $N$,
which proves the bound for $h$. The numerator of $h$
anticommutes with $Q$ to give $N$, squares to zero, and
commutes with $N$. The resulting polynomial identities
therefore extend by density and continuity.

The operator never differentiates a charge. Its charge
linearity extends from polynomials to entire functions.
\end{proof}

In particular, the inclusion from polynomial to analytic
coefficients changes the cohomology from
$\mathbb C[J_i:i\in I]$ to $\mathcal O(\mathbb C^I_J)$.
It is not a quasi-isomorphism. The larger coefficient algebra
is the cost of allowing arbitrary numerical charges and
convergent mode series.

\begin{theorem}[Translation and its optimal radius]
\label{cbc:ac:translation}
For $|t|<R$, translation extends continuously to
\[
 \tau_t=e^{tT}:\mathcal Y(R)\longrightarrow
                                      \mathcal Y(R-|t|).
\]
If $0<|t|<R$ and $S>R-|t|$, no continuous linear map
$\mathcal Y(R)\to\mathcal Y(S)$ agrees with $\tau_t$
on all mode polynomials.

For $0<a<R$ and $D_a=\{t:|t|<a\}$, the formula
\[
 E_t=(1+dt\,L)\tau_t
\]
defines a continuous dg algebra map
\[
 E:\mathcal Y(R)\longrightarrow
                         \Omega^\bullet(D_a;\mathcal Y(R-a)).
\]
It satisfies $E_{t+s}=E_tE_s$ wherever both sides have a
common radius domain.
\end{theorem}
\begin{proof}
For each generator,
\[
 \tau_tB_n=\sum_{j=0}^n\binom njt^{n-j}B_j,\qquad
 \|\tau_tB_n\|_{A,r}\leq A(r+|t|)^n.
\]
For $C_n$ the corresponding bound is
$Ar(r+|t|)^n\leq A(r+|t|)^{n+1}$.
Multiplicativity and submultiplicativity give
\begin{equation}\label{cbc:ac:translation-bound}
 \|\tau_tF\|_{A,r}\leq\|F\|_{A,r+|t|}.
\end{equation}
Every target radius satisfies $r+|t|<R$, so the
polynomial map extends continuously.

To prove optimality, write $t=s e^{i\alpha}$ with
$0<s<R$ and take
\[
 F=\sum_{n\geq0}R^{-n}e^{-in\alpha}B_n.
\]
Its seminorm at $A,r$ is $A/(1-r/R)$, so it belongs
to $\mathcal Y(R)$. In the translated polynomial
truncations, the coefficient of $B_j$ tends to
\[
 \sum_{n\geq j}R^{-n}e^{-in\alpha}\binom njt^{n-j}
       =e^{-ij\alpha}\frac{R}{(R-s)^{j+1}}.
\]

Coefficient evaluation is continuous. Any continuous
extension must therefore have these coefficients on $F$.
Their absolute sum diverges at every radius $r>R-s$.
Choosing $R-s<r<S$ proves the assertion.

For a compact $K\subset D_a$ and $r<R-a$, one has
$r+\sup_K|t|<R$. Thus \eqref{cbc:ac:translation-bound}
is uniform on $K$. The estimate for $L$ controls the
one-form component. On polynomials, the identities
$QL+LQ=T$, $[T,L]=0$, and $L^2=0$ give
\[
 (Q+d_t)E_t=E_tQ,\qquad E_{t+s}=E_tE_s.
\]

The operator $dt\,L$ is a degree-zero derivation whose
square is zero, so $1+dt\,L$ is multiplicative.
Continuity extends these identities to the stated spaces.
For two numerical translations a common target radius is
any positive radius below $R-|t|-|s|$.
\end{proof}

The primitive coproduct $\Delta_0$ is continuous for the
independently completed sites. For a target seminorm at
$(A_1,A_2;r_1,r_2)$, choose a source amplitude at least
$A_1+A_2$ and a source radius at least
$\max(r_1,r_2)$, still below $R$. These choices bound
the images of every generator. A translation in the first
factor replaces $r_1$ by $r_1+|t|$.
The spectral coproduct identities of
Proposition~\ref{cbc:cf:spectral} consequently extend on
the common domains given by these inequalities.

\section{The separated gauge and its collision boundary}
\label{cbc:ac:gauge}

For two sites with separation $z$, the full primitive of
\eqref{cbc:cf:primitive} and its differential are
\begin{align}
 b(z)&=-\frac1\kappa\sum_{N\geq1}\frac{z^{-N}}N
              \sum_{a+b=N}(-1)^a\binom Na
                         B_{1,a}B_{2,b},
                  \label{cbc:ac:b}\\
 r(z)&=\sum_{a,b\geq0}(-1)^a\binom{a+b}{a}
        \frac{C_{1,a}B_{2,b}-B_{1,a}C_{2,b}}
                    {z^{a+b+1}}.
                  \label{cbc:ac:r}
\end{align}
Let $L_1$ act at the first site, and put
\[
 D_0=Q+d_z,\qquad
 \mathcal A=r+dz\,L_1r,\qquad
 \rho=\kappa^{-1}J_1J_2\,\frac{dz}{z}.
\]
The formulas $Qb=r$ and
$d_zb=dz\,L_1r-\iota\rho$ hold coefficientwise.
We now determine their domain of convergence.

\begin{theorem}[Exact gauge domain]\label{cbc:ac:exact-domain}
Fix $R_1,R_2>0$ and $z\ne0$. Each of
$b(z)$, $e^{-b(z)}$, and $e^{b(z)}$ belongs to
$\mathcal Y_{\{1,2\}}(R_1,R_2)$ exactly when
\[
 |z|\geq R_1+R_2.
\]
Their largest open domain of holomorphic dependence in
this coefficient space is $|z|>R_1+R_2$.
For $q=(r_1+r_2)/|z|<1$,
\begin{equation}\label{cbc:ac:exact-norm}
 \|b(z)\|_{\mathbf A,\mathbf r}
   =\frac{A_1A_2}{|\kappa|}
                      \sum_{N\geq1}\frac{q^N}{N}
   =\frac{A_1A_2}{|\kappa|}\log\frac1{1-q}.
\end{equation}
The tail after weight $N_0$ is bounded by
\[
 \frac{A_1A_2}{|\kappa|}
       \frac{q^{N_0+1}}{(N_0+1)(1-q)}.
\]
\end{theorem}
\begin{proof}
The monomials $B_{1,a}B_{2,b}$ are distinct. Summing
their absolute coefficients at $a+b=N$ gives
$A_1A_2q^N/(|\kappa|N)$ by the binomial theorem.
The harmonic power series converges precisely for $q<1$.
Requiring this for every $r_i<R_i$ is equivalent to
$|z|\geq R_1+R_2$. The tail estimate follows by replacing
$1/N$ with $1/(N_0+1)$ in its geometric majorant.

Submultiplicativity gives
$\|e^{\pm b}\|\leq e^{\|b\|}$ wherever $b$ belongs to
the coefficient space. Conversely, projection to total
field degree two is continuous and sends $e^{\pm b}$ to
$\pm b$. The exponentials are defined coefficientwise,
since $b$ has positive jet weight. Their membership
therefore implies that of $b$.

All bounds are uniform on compact subsets of the
strict-inequality domain. Differentiated series converge
there by the same geometric estimates, which prove
holomorphic dependence. Every neighborhood of a boundary
point contains positions where membership fails.
Boundary membership supplies no larger open domain.
\end{proof}

\begin{theorem}[Analytic module contraction]
\label{cbc:ac:gauge-retract}
Let $U\subset\{z:|z|>R_1+R_2\}$ be open and put
$G=e^{-b}$. There is a continuous strong deformation
retract of module complexes
\[
 \bigl(\Omega^\bullet(U;\mathcal Y_{\{1,2\}}(R_1,R_2)),
                                      D_0+\mathcal A\cdot\bigr)
 \ \rightleftarrows\
 \bigl(\Omega^\bullet(U;\mathcal O(\mathbb C^2_J)),
                                      d_z+\rho\cdot\bigr)
\]
with inclusion $G\iota$, projection $\pi$, and homotopy
$GhG^{-1}$. At any defining seminorm,
\[
 \|GhG^{-1}\|\leq|\kappa|^{-1}e^{2\|b\|}.
\]
\end{theorem}
\begin{proof}
In addition to \eqref{cbc:ac:exact-norm}, the binomial
sums in \eqref{cbc:ac:r} give
\[
 \|r(z)\|\leq
      \frac{A_1A_2(r_1+r_2)}{|z|-r_1-r_2},
 \qquad
 \|L_1r(z)\|\leq
      \frac{A_1A_2}{|\kappa|(|z|-r_1-r_2)}.
\]
They make the connection form and the coefficient
identities convergent on $U$.

To check the first identity
directly, differentiating a first-site boson in
\eqref{cbc:ac:b} uses
\[
 -\frac{(-1)^a a}{N}\binom Na
                 =(-1)^{a-1}\binom{N-1}{a-1}.
\]
Differentiating a second-site boson supplies the other
term of $r$. Position differentiation of $-z^{-N}/N$
gives $z^{-N-1}dz$. Its missing weight-zero term is
$\iota\rho$. Hence
\[
 D_0b=\mathcal A-\iota\rho,\qquad
 (D_0+\mathcal A\cdot)G=G(D_0+\iota\rho\cdot).
\]

Extend $h$ over forms with the Koszul sign. Since it is
charge-linear, it anticommutes with both $d_z$ and
multiplication by $\rho$. Proposition
~\ref{cbc:ac:continuous-contraction} thus contracts the
$\rho$-twisted complex onto the charge connection.
Conjugation by $G$ gives the asserted retract.

Every monomial of $b$ contains a positive-index boson,
so $\pi G=\pi G^{-1}=\pi$. This proves that the
projection remains $\pi$ and that $\pi G\iota=1$.
The homotopy bound follows by multiplying the three
operator bounds.
\end{proof}

For a finite ordered set of positions, require
$|z_i-z_j|>R_i+R_j$ for every pair and define
\[
 b_I=\sum_{i<j}b_{ij}(z_i-z_j),\qquad
 \mathcal A_I=\sum_{i<j}(r_{ij}+dz_{ij}L_i r_{ij}),
 \qquad
 \rho_I=\kappa^{-1}\sum_{i<j}J_iJ_j\,d\log z_{ij}.
\]
All even primitives commute. The estimates add over
finitely many pairs, so the same retract holds with
$G_I=e^{-b_I}$. Every pair is retained, including the
nonconsecutive channel at three positions and the
crossed channels at four.

If a specified continuous coefficient map $f$ relates
two untwisted models on a common radius domain, its
conjugate is $G_t^{-1}fG_s$. For multiplicative $f$ its
extra factor is $\exp(b_t-f(b_s))$. In compositions the
intermediate factors cancel. Thus conjugation preserves
the identities of the supplied maps. It does not supply
a missing map from a native collision coefficient.
Also $G\iota(1)=e^{-b}$: the endpoint of this module
contraction is not the constant unit.

\begin{proposition}\label{cbc:ac:essential}
The full gauge need not extend through $z=0$ as a
finite-pole meromorphic coefficient map on the same
analytic carrier.
\end{proposition}
\begin{proof}
Set all modes equal to zero except $B_{1,1}=a$ and
$B_{2,0}=q$, with $aq\ne0$. This coefficient evaluation
is continuous. For any fixed $r_1>0$, choose
$A_1r_1>|a|$ and $A_2>|q|$.
It sends
\[
 b(z)\longmapsto\frac{aq}{\kappa z},\qquad
 G(z)\longmapsto\exp\left(-\frac{aq}{\kappa z}\right).
\]

The latter has nonzero coefficients in every negative
power of $z$. No power of $z$ makes it holomorphic.
A finite-pole extension would remain finite-pole under
the continuous evaluation, which is impossible.
\end{proof}

The logarithmic connection has a different obstruction:
its cohomology changes at resonant charges.

\begin{proposition}\label{cbc:ac:resonance}
At numerical charges $q_1,q_2$, put $\nu=q_1q_2/\kappa$.
On a holomorphic annulus, the complex
$d_z+\nu\,dz/z$ is acyclic if $\nu\notin\mathbb Z$.
If $\nu\in\mathbb Z$, its cohomology is one-dimensional
in each of degrees zero and one, represented by
$z^{-\nu}$ and $z^{-\nu}dz/z$.
Its horizontal monodromy is $e^{-2\pi i\nu}$.
There is no contraction over the entire charge algebra
that eliminates every nonzero Laurent exponent uniformly
in the charges.
\end{proposition}
\begin{proof}
On the Laurent mode $z^m$, the differential is
multiplication by $m+\nu$, with target $z^m dz/z$.
Except at the possible resonant exponent, division by
$m+\nu$ has at most polynomial growth in $m$. It
preserves Laurent convergence on compact subannuli
and contracts exactly those modes. The horizontal
solution is $z^{-\nu}$ on a simply connected cover,
which gives the monodromy.

For any nonzero integer $m$, the entire-charge inverse
$(m+J_1J_2/\kappa)^{-1}$ has a pole along the nonempty
hypersurface $J_1J_2=-\kappa m$. At a charge on this
hypersurface the mode has nonzero cohomology.
Uniform elimination would contradict that specialization.
An inverse exists as a formal charge series near zero,
but it does not extend to an entire charge function.
\end{proof}

The essential singularity concerns the full-mode gauge.
Resonance concerns the residual charge connection.
A collision comparison must account for both on its
declared coefficient category. Finite algebraic models
retain the regular attachment through a collision.

\section{Punctures with their regular coefficient}
\label{cbc:ac:punctures}

Let $k$ be a field of characteristic zero, and choose distinct
$a_1,\ldots,a_r\in k$, with $r\geq1$. Put
\[
 R=k[t],\qquad f=\prod_{i=1}^r(t-a_i),\qquad
 L_f=R[f^{-1}],\qquad
 D=k\langle t,\partial\rangle/(\partial t-t\partial-1).
\]
All differential modules in this section are left $D$-modules.
The regular and localized modules have their derivative
connections. For each puncture, the delta module is
\[
 \delta_i=D/D(t-a_i)
              \simeq R[(t-a_i)^{-1}]/R.
\]
Partial fractions give the exact sequence
\begin{equation}\label{cbc:ac:puncture-extension}
 0\longrightarrow R\longrightarrow L_f
                 \longrightarrow\bigoplus_{i=1}^r\delta_i
                 \longrightarrow0.
\end{equation}

We will compare two choices of objects in the dg category
generated by this sequence. Keeping objects \emph{marked}
means requiring equivalences to preserve the specified
objects individually.

For $q\in D$, write $r_q$ for right multiplication by $q$.
The free resolutions
\begin{equation}\label{cbc:ac:normal-resolutions}
 P_0=[D\xrightarrow{r_\partial}D]\longrightarrow R,\qquad
 P_i=[D\xrightarrow{r_{t-a_i}}D]\longrightarrow\delta_i
\end{equation}
have terms in degrees $-1,0$. Their differentials are
injective: the order filtration identifies the associated
graded of $D$ with the polynomial domain $k[t,\xi]$.
The degree-zero augmentations send $1$ to the cyclic
generator of the target. Products of maps below are
compositions, with the right map applied first.

\begin{lemma}\label{cbc:ac:puncture-hom}
The mapping complexes of the regular and delta modules
have cohomology
\[
 \begin{array}{c|cc}
  \operatorname{RHom}_D(-,-)&R&\delta_j\\ \hline
  R&k&k[-1]\\
  \delta_i&k[-1]&
       \begin{cases}k,&i=j,\\0,&i\ne j\end{cases}
 \end{array}.
\]
Here $k[-1]$ is in cohomological degree one.
\end{lemma}
\begin{proof}
Mapping the two resolutions into $R$ gives the operators
$\partial$ and $t-a_i$ on polynomials. Integration
contracts the first to constants. The second is injective
with one-dimensional cokernel.

On $\delta_j$, differentiation is injective with cokernel
the simple pole $(t-a_j)^{-1}$. Multiplication by $t-a_j$
is surjective, with kernel that same pole. If $i\ne j$,
then $t-a_i=(t-a_j)+(a_j-a_i)$ is invertible on
$\delta_j$. Its inverse on each vector is a finite
geometric sum because $t-a_j$ acts locally nilpotently.
These calculations give all four kinds of corners.
\end{proof}

Define a finite graded algebra $A_r$ with identities
$e_0,\ldots,e_r$ and degree-one arrows
\[
 \alpha_i:0\longrightarrow i,\qquad
 \beta_i:i\longrightarrow0\quad(1\leq i\leq r).
\]
Every product of two arrows is zero, and the differential
is zero. Its dimension is $3r+1$. The central vertex
represents the regular module and the leaves represent
the individual closed coefficients.

Write $\operatorname{Perf}(A_r)$ for perfect right dg
$A_r$-modules. Write $\operatorname{thick}_D$ for closure
under finite cones, shifts, and direct summands in the
derived category of left $D$-modules.

\begin{theorem}[The formal star category]
\label{cbc:ac:star}
There is a strict dg-algebra quasi-isomorphism
\[
 A_r\longrightarrow
           \operatorname{End}_D(P_0\oplus P_1\oplus\cdots\oplus P_r).
\]
The dg Yoneda functor on these resolutions gives
\[
 \operatorname{Perf}(A_r)\simeq
 \operatorname{thick}_D(R,\delta_1,\ldots,\delta_r)
                  =\operatorname{thick}_D(R\oplus L_f).
\]
\end{theorem}
\begin{proof}
Send $e_i$ to the projection onto $P_i$. Send $\alpha_i$
to the degree-one identity $P_0^{-1}\to P_i^0$ and
$\beta_i$ to the degree-one identity $P_i^{-1}\to P_0^0$.
These maps are closed. Their pairwise products vanish
strictly, since their targets have degree zero, on which
the next degree-one map vanishes. They represent the
nonzero cokernel generators in
Lemma~\ref{cbc:ac:puncture-hom}. The map is therefore
a quasi-isomorphism on every corner.

The dg Yoneda functor identifies finite cones and
summands of these resolutions with perfect $A_r$-modules.
Sequence~\eqref{cbc:ac:puncture-extension} puts $L_f$
in the category generated by the regular and delta
modules. Conversely, its cone recovers the direct sum
of the deltas, and direct summands recover each delta.
\end{proof}

Each component of the extension
\eqref{cbc:ac:puncture-extension} is nonsplit.
A splitting would embed a nonzero vector annihilated
by $t-a_i$ into the torsion-free module $L_f$.
Normalize its nonzero extension class to one.
Applying derived Hom into $R$ then gives
\begin{equation}\label{cbc:ac:reverse-localization}
 \operatorname{RHom}_D(L_f,R)
       \simeq[k\xrightarrow{(1,\ldots,1)}k^r]
       \simeq k^{r-1}[-1].
\end{equation}

The two-term complex has degrees zero and one.
The reverse Hom vanishes for one puncture and acquires
a degree-one class at the second puncture.
Its compositions will detect the localization attachment.

\subsection{The marked pair and its ternary products}

The twisted complexes needed here are finite sums of
objects of $A_r$, with a strictly triangular degree-one
matrix $d$ satisfying $d^2=0$. Define
\[
 O=e_0,\qquad
 \mathcal L=e_0\oplus e_1\oplus\cdots\oplus e_r,\qquad
 d_{\mathcal L}=\sum_i\beta_i.
\]
With $e_0$ first, the twisting matrix is triangular.

The corresponding free resolution has generators
$v_0,\ldots,v_r$ in degree zero and relations
\[
 \partial v_0=0,\qquad v_0+(t-a_i)v_i=0.
\]
Its augmentation sends $v_0$ to $1$ and $v_i$ to
$-(t-a_i)^{-1}$. This identifies its cokernel with $L_f$:
the regular submodule and each quotient delta have the
extension components just normalized.

The differential
is injective. A vector in its kernel first satisfies
$x_i(t-a_i)=0$ at every leaf, hence $x_i=0$.
Its remaining equation is $x_0\partial=0$, hence $x_0=0$.

Set $B_r=\operatorname{End}(O\oplus\mathcal L)$ in the
category of twisted complexes. Before twisting, write
$\mathcal L_j$ for its $j$-th summand. A basis is
\[
 \begin{array}{c|c|c}
  \text{element}&\text{matrix entry}&\text{degree}\\ \hline
  e&O\to O&0\\
  a&O\to\mathcal L_0&0\\
  p&\mathcal L_0\to O&0\\
  u_j&\mathcal L_j\to\mathcal L_j\ (0\leq j\leq r)&0\\
  b_i&O\to\mathcal L_i\ (\alpha_i)&1\\
  c_i&\mathcal L_i\to O\ (\beta_i)&1\\
  \theta_i&\mathcal L_0\to\mathcal L_i\ (\alpha_i)&1\\
  z_i&\mathcal L_i\to\mathcal L_0\ (\beta_i)&1 .
 \end{array}
\]

Multiply matching entries by matrix composition. Every
product of two degree-one entries is zero.
The commutator differential with the twisting matrix is
\begin{equation}\label{cbc:ac:marked-differential}
 dp=-\sum_i c_i,\qquad
 du_0=-\sum_i z_i,\qquad du_i=z_i\quad(i>0),
\end{equation}
and is zero on the remaining basis. In particular,
\[
 pa=e,\qquad ap=u_0,\qquad
 ac_i=z_i,\qquad \theta_i a=b_i,\qquad b_i p=\theta_i.
\]
This specifies a dg algebra of dimension $5r+4$.

Its cohomology $H$ has degree-zero representatives
$E=e$, $F=\sum_{j=0}^r u_j$, and $a$. In degree one,
choose $b_i,\theta_i$ for $i\leq r$ and $c_j$ for
$j<r$. The remaining class is
$[c_r]=-\sum_{j<r}c_j$. Thus $\dim H=3r+2$.
Apart from the two identities, the only nonzero binary
products are $\theta_i a=b_i$.

Let $i:H\to B_r$ be the chosen inclusion. Define
$\pi:B_r\to H$ by
\[
 \pi(u_0)=F,\quad \pi(u_j)=0\ (j>0),\quad
 \pi(p)=\pi(z_i)=0,\quad
 \pi(c_r)=-\sum_{j<r}c_j,
\]
and by the identity on the other chosen representatives.
Define a degree-minus-one map $h$ by
\[
 h(z_i)=u_i,\qquad h(c_r)=-p,
\]
and set it to zero on the other basis vectors.
Equation~\eqref{cbc:ac:marked-differential} gives
\[
 \pi i=1,\qquad dh+hd=1-i\pi,\qquad
 h^2=hi=\pi h=0.
\]
These formulas retain the homotopies of the products
$ac_j$, which vanish in cohomology.

We use the suspension $|sx|=|x|-1$ and the conventions
\[
 b_2(sx,sy)=(-1)^{|x|}s(xy),\qquad
 b_3(sx,sy,sz)=(-1)^{|y|}s\,m_3(x,y,z).
\]
Thus an unsuspended $m_n$ has degree $2-n$.

\begin{theorem}\label{cbc:ac:marked-model}
The algebra $B_r$ models
$\operatorname{REnd}_D(R\oplus L_f)$ with its two
objects marked. Its strictly unital minimal
$A_\infty$-model on $H$ has the binary product above
and the following nonzero ternary products on
nonidentity basis elements:
\begin{align}
 m_3(a,c_j,\theta_i)&=-\delta_{ij}\theta_i,&
 m_3(a,c_j,b_i)&=-\delta_{ij}b_i,\notag\\
 m_3(c_\ell,a,c_j)&=-\delta_{\ell j}c_\ell
        &&(j,\ell<r,\ i\leq r).
                  \label{cbc:ac:ternary}
\end{align}
All $m_n$ with $n\geq4$ vanish. A suspended comparison
to $B_r$ has $F_1=i$ and
$F_2(sa,sc_j)=su_j$ for $j<r$.
Its other higher components vanish.
\end{theorem}
\begin{proof}
The strict quasi-isomorphism of
Theorem~\ref{cbc:ac:star} extends to finite triangular
twisted complexes and their Hom complexes. The displayed
resolution identifies $O,\mathcal L$ with $R,L_f$,
including their compositions.

Transfer through $(i,\pi,h)$ gives
\[
 m_3(x,y,z)=(-1)^{|x|}
       \pi\bigl(i(x)h(i(y)i(z))\bigr)
       -\pi\bigl(h(i(x)i(y))i(z)\bigr).
\]
The only product of two representatives on which $h$
is nonzero is $ac_j=z_j$, with $j<r$, and its
homotopy is $u_j$. The products
$u_j\theta_i$, $u_jb_i$, and $c_\ell u_j$
give exactly \eqref{cbc:ac:ternary}.

Multiplying $u_j$ by a representative gives zero,
$u_j$, or a representative killed by $h$. A transfer
tree with another homotopy above this multiplication
therefore vanishes. The remaining four-input tree has
root product $u_i u_j$, whose projection is zero.
Every larger tree contains one of these configurations.
The same reasoning leaves only the stated $F_2$.
Since $hi=0$, the two identities remain strict units.

The bar equation $b_{B_r}F=Fb_H$ can also be checked
without a tree expansion. Its nonzero terms reduce
to $du_j=z_j$ when one pair $(a,c_j)$ occurs,
and to $u_i u_j=\delta_{ij}u_j$ when two such
pairs occur. The suspended signs are those fixed
above. All other terms vanish by the matrix products.
This verifies the comparison in every arity.
\end{proof}

\begin{theorem}\label{cbc:ac:marked-nonformality}
For $r\geq2$, the marked dg category $(R,L_f)$ is
not equivalent, by an object-preserving strictly
unital $A_\infty$-equivalence, to its binary
cohomology category.
\end{theorem}
\begin{proof}
Normalize the linear term of a hypothetical equivalence
by the inverse of its cohomology isomorphism. Its
linear term is then the identity. The arity-three
equation would express $m_3$ as the Hochschild boundary
of a binary map $G_2$ of degree minus one.

Evaluate that boundary on $(a,c_j,\theta_i)$.
Both adjacent binary products are zero. The term
$aG_2(c_j,\theta_i)$ also vanishes: its second
factor is a reverse degree-one map, and $ac_\ell=0$
in $H$. Only $G_2(a,c_j)\theta_i$ can remain,
up to its fixed sign. But
$G_2(a,c_j)\in H^0\operatorname{End}(\mathcal L)=kF$,
so it acts by the same scalar on every $\theta_i$.

For fixed $j<r$, formula~\eqref{cbc:ac:ternary}
acts by $-1$ on $\theta_j$ and by zero on
$\theta_r$. No scalar has both values.
The resulting rank-one projector survives modulo
scalar endomorphisms, a contradiction.
\end{proof}

For two punctures, the obstruction is already
\[
 m_3(a,c_1,\theta_1)=-\theta_1,\qquad
 m_3(a,c_1,\theta_2)=0.
\]
For one puncture, there is no reverse class and
the five-dimensional minimal model is formal.
The formal star generator and the nonformal marked
pair can therefore generate the same dg category.
The marked objects determine which formality question
is being asked.

\subsection{The primitive three-point chart}

The primitive module $H_3$ of
\eqref{cbc:hc:primitive} supplies a geometric instance.
On $z_1-z_3\ne0$, write
\[
 z_3=s_0,\qquad z_1=s_0+u,\qquad z_2=s_0+ut.
\]
Here $u$ is invertible. Let $e_i$ denote the generic
residue vectors, and set $v_0=e_2-e_3$ and
$v_1=e_1-e_2$.

\begin{proposition}\label{cbc:ac:primitive-chart}
There is a horizontal isomorphism
\[
 H_3|_{u\ne0}\simeq k[s_0,u,u^{-1}]\otimes_k
 \bigl(k[t]e_3\oplus k[t,t^{-1}]v_0
                   \oplus k[t,(t-1)^{-1}]v_1\bigr).
\]
Its generated dg category has model
$\operatorname{Perf}(E\otimes_k A_2)$, where
$E=\Lambda(\rho)$, $|\rho|=1$, and $d\rho=0$.
Adding the auxiliary exterior coefficients of $H_3$
does not enlarge this generated category.
\end{proposition}
\begin{proof}
The finite presentation of
Theorem~\ref{cbc:hc:presentation} gives
\[
 g_{23}=\frac{v_0}{ut},\qquad
 g_{12}=\frac{v_1}{u(1-t)},\qquad
 g_{13}=\frac{v_0+v_1}{u},\qquad
 g_{123}=-\frac{v_0}{u^2t}+\frac{v_1}{u^2(1-t)}.
\]
The singleton generators are regular combinations of
$e_3,v_0,v_1$. These formulas and their position
derivatives place every generator in the proposed
module. Conversely, $g_{23},g_{12}$ and their
derivatives generate all finite poles at $t=0,1$
in the two stated lines.

Flatness from
Corollary~\ref{cbc:hc:flat} implies torsion-freeness.
The generic residue map is thus injective, and the
two submodules with these generators coincide.
Differentiating the rational representatives proves
horizontality.

The three shape summands generate $k[t],\delta_0,\delta_1$
by their localization triangles. Theorem~\ref{cbc:ac:star}
identifies their category with $\operatorname{Perf}(A_2)$.
It also contains the double localization. Exterior
coefficients are sums of these localizations. A rank-one
localized summand has zero second exterior power,
as follows by flatness and passage to the generic point.
No additional generator is needed.

The scale coefficient $k[u,u^{-1}]$ has de Rham
cohomology represented by $1,du/u$. Their span is
the strict algebra $E$. Polynomial common translation
contracts to $k$. Tensoring the bounded free
resolutions in the three independent coordinates
gives the asserted dg model.
\end{proof}

The three summands of the primitive are a different
marking from the pair consisting of a regular and
a doubly localized module.

\begin{proposition}\label{cbc:ac:primitive-formal}
Put $L_1=k[t,t^{-1}]$, $L_2=k[t,(t-1)^{-1}]$,
and $K=k[t]\oplus L_1\oplus L_2$, with the three
summands marked. Its derived endomorphism algebra
has a strict formal model $C$ of dimension eleven.
The degree-zero basis consists of the three identities
and $a_i:O\to L_i$ for $i=1,2$. The degree-one
basis consists of
\[
 b_i:O\to L_i,\qquad
 \theta_{ji}:L_i\to L_j\quad(i,j=1,2).
\]
Apart from identity actions, the products are
$\theta_{ji}a_i=b_j$. The marked primitive on
the chart consequently has formal model $E\otimes C$.
\end{proposition}
\begin{proof}
Over $A_2$, use $O=e_0$ and the twisted complexes
$\mathcal L_i=(e_0\oplus e_i,\beta_i)$.
Represent $a_i$ by the identity into the central
summand and $b_i$ by $\alpha_i$ into the leaf.
Represent $\theta_{ji}$ by $\alpha_j$ from the central
summand of $\mathcal L_i$ to the leaf of
$\mathcal L_j$. All these maps are closed.

Each reverse block $\operatorname{Hom}(\mathcal L_i,O)$
is a two-term isomorphism and is acyclic. For $i\ne j$,
the block $\operatorname{Hom}(\mathcal L_i,\mathcal L_j)$
has one central map in degree zero and two maps in
degree one. Its differential spans the map from the
source leaf to the target center, leaving $\theta_{ji}$.
On a diagonal block, the sum of its two identities
and $\theta_{ii}$ represent the two classes.
The remaining blocks are represented by $1_O,a_i,b_i$.

These eleven maps therefore represent every class and
no boundary. They form a subalgebra with the stated
products, so their inclusion is a strict
quasi-isomorphism. Tensoring with the scale model and
the contracted translation factor proves the last claim.
\end{proof}

The exterior coefficients $\Lambda^0H_3$ and
$\Lambda^3H_3$ restrict instead to the scale factor
tensored with $k[t]$ and
$k[t,t^{-1},(t-1)^{-1}]$. Marking this pair retains
the ternary projector of
Theorem~\ref{cbc:ac:marked-nonformality}.
Tensoring with $E$ does not remove it: the degree-zero
endomorphisms remain scalars, and the reverse
degree-one maps are still the $c_j$.
Thus primitive formality does not determine the
marked exterior system.

\section{Two punctures through a collision}
\label{cbc:ac:confluence}

We now allow the punctures to meet. Work over $B=k[s]$
with the relative Weyl algebra
\[
 D_B=B\langle t,\partial_t\rangle/
                         (\partial_t t-t\partial_t-1).
\]
The parameter $s$ is central. At this stage there is
no derivative in $s$. Put
\[
 f=t(t-s),\qquad L_B=B[t,f^{-1}],\qquad
 A_m=f\partial_t+m f_t,\qquad
 P_m=D_B/D_BA_m\quad(m\geq1).
\]

The cyclic generator $e_m$ maps to $f^{-m}$.
The relation $A_m(f^{-m})=0$ makes this a
differential-module map.

The finite-stage phenomenon of
Example~\ref{cbc:hc:spurious-delta} is explicit here.
For a nonzero scalar $a$, the specialization $s=a$
has two simple roots, and $P_m$ presents their
localization.

Near a root, write $f=uv$ with $u$
its linear coordinate and $v$ invertible. Replacing
the generator by $v^m e_m$ gives the relation
$(u\partial_u+m)v^m e_m=0$. The Weyl normal
form identifies this module with the localization
generated by $u^{-m}$. Away from both roots it
is the regular rank-one connection. These local
presentations establish the claim on the fiber.

At $s=0$, put $C_m=t\partial_t+2m$.
Then $A_m=tC_m$, and $g_m=C_me_m$ satisfies
\[
 tg_m=0,\qquad C_mg_m=(2m-1)g_m.
\]
There is an exact sequence
\[
 0\longrightarrow D/Dt
   \xrightarrow{\,r_{C_m}\,}D/D(tC_m)
   \longrightarrow D/DC_m\longrightarrow0.
\]

Right cancellation in $D$ proves injectivity of the
first map. The last quotient is $k[t,t^{-1}]$,
generated by $t^{-2m}$. The vector
$\ell_m=e_m-g_m/(2m-1)$ satisfies $C_m\ell_m=0$
and splits this quotient. Hence
\begin{equation}\label{cbc:ac:extra-delta}
 P_m\otimes_B k_0\simeq k[t,t^{-1}]\oplus\delta_0,
 \qquad g_m\longmapsto0\text{ in }L_B\otimes_B k_0.
\end{equation}
The finite presentation acquires an extra closed
coefficient at the collision.

The transition between denominator orders removes it.
The identity
\[
 A_mf=fA_{m+1}
\]
gives $P_m\to P_{m+1}$ by $e_m\mapsto f e_{m+1}$.
It lifts to right multiplication by $f$ in both
degrees of $[D_B\xrightarrow{r_{A_m}}D_B]$.
Lemma~\ref{cbc:hc:localization-telescope} identifies
the direct-sum telescope of these resolutions with
$L_B$.

Explicitly, the colimit relations give
\[
 \partial_t e_m=-m f_t e_{m+1},\qquad
 \frac p{f^m}\longleftrightarrow p e_m.
\]
These maps are inverse and horizontal.
They remain so after every scalar extension $B\to C$:
the leading coefficient $f$ is monic in $t$, hence
a non-zero-divisor over $C$.

At $s=0$, the transition sends
\[
 g_m\longmapsto C_m t^2e_{m+1}
                 =t^2C_{m+1}e_{m+1}=t^2g_{m+1}=0.
\]
Thus the full telescope retains exactly the
localization at the collision.

\subsection{Finite traces and the inverse telescope}

Apply derived Hom into $B[t]$ with its derivative
connection. The stage complex is
$[B[t]\xrightarrow{A_m}B[t]]$ in degrees zero and one.
On monomials,
\begin{equation}\label{cbc:ac:stage-action}
 A_m(t^n)=(n+2m)t^{n+1}-(n+m)s t^n.
\end{equation}
Write $(a)_n=a(a+1)\cdots(a+n-1)$, with $(a)_0=1$.

\begin{lemma}\label{cbc:ac:finite-trace}
The stage complex has zero kernel and a free rank-one
cokernel over $B$. Its quotient map, normalized by
$\pi_m(1)=1$, is
\[
 \pi_m(t^n)=\frac{(m)_n}{(2m)_n}s^n.
\]
The inverse transition on cokernels satisfies
\begin{equation}\label{cbc:ac:trace-transition}
 \pi_m(fp)=-\frac{m}{2(2m+1)}s^2\pi_{m+1}(p).
\end{equation}
\end{lemma}
\begin{proof}
The leading coefficient in \eqref{cbc:ac:stage-action}
is an invertible positive integer. The operator
raises the degree of every nonzero polynomial by one,
so it is injective. Descending elimination reduces
each polynomial uniquely to a constant modulo its
image. The successive relations are
\[
 \pi_m(t^{n+1})
       =\frac{n+m}{n+2m}s\,\pi_m(t^n).
\]
They give the stated moments and a polynomial
primitive for every vector with zero moment.
Substituting $f=t^2-st$ and these moments proves
\eqref{cbc:ac:trace-transition} on each monomial.
\end{proof}

Rescale by constants $c_1=1$ and
$c_{m+1}=-m c_m/(2(2m+1))$.
In the coordinates $a_m=c_m\pi_m(p_m)$, the
inverse system becomes
\[
 B\xleftarrow{s^2}B\xleftarrow{s^2}B
                           \xleftarrow{s^2}\cdots
\]
in degree one. Dualizing the full direct-sum telescope
requires products and the derived inverse limit of
this system.

\begin{theorem}\label{cbc:ac:relative-dual}
With $Q=k[[s]]/k[s]$ taken as an algebraic module
quotient, there is an equivalence of $B$-complexes
\[
 \operatorname{RHom}_{D_B}(L_B,B[t])
                 \simeq Q[-2].
\]
In the normalized coordinates, the quotient is realized by
\[
 (a_m)_{m\geq1}\longmapsto
            \left[\sum_{m\geq1}s^{2m-2}a_m\right].
\]
\end{theorem}
\begin{proof}
Hom from the free telescope is the product difference
complex of the stage mapping complexes.
Lemma~\ref{cbc:ac:finite-trace} reduces it to
\[
 \prod_{j\geq0}B\xrightarrow{\Delta}\prod_{j\geq0}B,
 \qquad \Delta(b)_j=b_j-s^2b_{j+1},
\]
in degrees one and two. Its kernel is zero:
the relation forces each $b_j$ to be divisible
by every power of $s^2$.

The sum $\sum_{j\geq0}s^{2j}a_j$ is well-defined
as a formal series. On a boundary it telescopes
to the polynomial $b_0$. Conversely, suppose
the sum is a polynomial $p$. Set
\[
 b_j=s^{-2j}\left(p-\sum_{i<j}s^{2i}a_i\right).
\]

The numerator is a polynomial. Its formal-series
expression is a tail of order at least $2j$,
so it is divisible by $s^{2j}$ in $B$.
Thus $b_j\in B$ and $\Delta(b)=a$.
Finally, any $\sum c_n s^n$ is obtained by taking
$a_j=c_{2j}+c_{2j+1}s$.
This proves the asserted cokernel and all its degrees.
\end{proof}

The constant sequence $a_m=1$ gives the nonzero
class of $(1-s^2)^{-1}$ modulo polynomials.
The formal-series quotient therefore represents
an actual relative mapping class.

\begin{proposition}\label{cbc:ac:base-change}
The preceding derived Hom does not commute with the
flat extension $B\to B[s^{-1}]$. More precisely,
\begin{align*}
 \operatorname{RHom}_{D_B}(L_B,B[t])
                   \otimes_B B[s^{-1}]&\simeq Q[-2],\\
 \operatorname{RHom}_{D_{B[s^{-1}]}}
       (B[s^{-1}][t,f^{-1}],B[s^{-1}][t])
                              &\simeq B[s^{-1}][-1].
\end{align*}
For a closed fiber $s=a$, however,
\[
 Q[-2]\otimes_B^{\mathbf L}k_a
       \simeq
       \begin{cases}k[-1],&a\ne0,\\0,&a=0\end{cases}.
\]
These are the derived Hom groups of the corresponding
closed-fiber localizations.
\end{proposition}
\begin{proof}
Multiplication by $s$ is bijective on $Q$.
If $sF$ is polynomial, then $F$ is polynomial.
Removing a constant term and dividing by $s$ gives
surjectivity. Thus localizing $Q$ leaves it unchanged.
In contrast, after inverting $s$ in the stage system,
all transitions are invertible. Its limit is
$B[s^{-1}]$, its higher derived limit vanishes,
and its surviving degree is one.

For the closed fibers use
$0\to B\to k[[s]]\to Q\to0$.
The formal-series module is torsion-free over the
principal ideal domain $B$, hence flat.
For $a\ne0$, the element $s-a$ is invertible on
$k[[s]]$. The exact triangle therefore gives
$Q\otimes_B^{\mathbf L}k_a\simeq k[1]$.
For $a=0$, bijectivity of $s$ on $Q$ makes the
tensor complex acyclic. Formula
~\eqref{cbc:ac:reverse-localization} gives the
same answers for two distinct punctures and one
puncture, respectively.
\end{proof}

In particular, $L_B$ is not perfect over $D_B$.
A bounded resolution by finite projectives would
make Hom into $B[t]$ commute with this flat extension.
The bounded free telescope has infinitely generated
terms, and its Hom involves the products just computed.

\subsection{Restoring the parameter connection}

The quotient $Q$ was computed without differentiating
in $s$. The flat telescope connection of
\eqref{cbc:hc:spatial-connection} restores that
derivative before taking Hom. On its lowest
generator it is
$\nabla_s h_m=-m f_s h_{m+1}$.
For the corresponding degree-two Hom coordinate
$p_m$, the induced action is therefore
\[
 p_m\longmapsto \partial_s p_m+m f_s p_{m+1}.
\]
This determines the connection on the quotient.

\begin{theorem}\label{cbc:ac:quotient-connection}
On $Q=k[[s]]/k[s]$, the induced parameter derivative is
\[
 D_Q[F]=\left[F'+\frac{F-F(0)}s\right].
\]
It satisfies $[D_Q,s]=1$ and is invertible as a
$k$-linear map. Its inverse is
\[
 K_Q[F]=\left[\frac1s\int_0^s\tau F(\tau)\,d\tau\right].
\]
Consequently, for the full Weyl algebra $D_{s,t}$,
\[
 \operatorname{RHom}_{D_{s,t}}(L_B,k[s,t])\simeq0.
\]
\end{theorem}
\begin{proof}
Choose constant-in-$t$ representatives $p_m=b_m(s)$
of the stage cokernels. Since $f_s=-t$ and
$\pi_m(t)=s/2$, their derivatives have moments
$b_m'-(m/2)s b_{m+1}$.
Multiplication by $c_m$ converts this to
\[
 a_m'+(2m+1)s a_{m+1}.
\]
For $F=\sum_{m\geq1}s^{2m-2}a_m$, summing gives
$F'+(F-a_1)/s$. Its difference from the stated
formula is the polynomial $(F(0)-a_1)/s$:
$a_1$ is a polynomial with constant term $F(0)$.
Thus both define the same map on $Q$.

Both $D_Q$ and $K_Q$ preserve polynomial
representatives. On formal monomials,
$K_Q(s^n)=s^{n+1}/(n+2)$.
Direct calculation gives
$D_QK_QF=F$ and $K_QD_QF=F-F(0)$.
The latter constant disappears in $Q$.
Also $D_Q(sF)-sD_Q(F)=F+F(0)$ on representatives,
which proves the Weyl relation on the quotient.

The Spencer resolution of
Proposition~\ref{cbc:hc:spencer-resolution} computes
full spatial Hom from the relative complex with
its parameter derivative. The relative cohomology
is concentrated in degree two, so truncation leaves
$[Q\xrightarrow{D_Q}Q][-2]$.
Its inverse $K_Q$ contracts the complex.
\end{proof}

For example, the relative class of $(1-s^2)^{-1}$
has the spatial primitive
\[
 -\frac{\log(1-s^2)}{2s}
     =\sum_{n\geq0}\frac{s^{2n+1}}{2n+2}.
\]
This is a formal series regular at zero.
No analytic convergence is needed in this computation.

There is also a finite verification of the full
spatial vanishing. Put $x=t$, $y=t-s$.
Then $L_B=k[x^{\pm1},y^{\pm1}]$. The two commuting
cyclic operators are $E_x=x\partial_x+1$ and
$E_y=y\partial_y+1$. Their tensor resolution gives
the Hom complex
\[
 k[x,y]\xrightarrow{p\mapsto(E_xp,E_yp)}k[x,y]^2
       \xrightarrow{(a,b)\mapsto E_xb-E_ya}k[x,y].
\]

Since $E_x^{-1}(x^iy^j)=x^iy^j/(i+1)$, its
contraction is
$h_1(a,b)=E_x^{-1}a$ and
$h_2(c)=(0,E_x^{-1}c)$.
Commutation with $E_y$ proves $dh+hd=1$ and
$h^2=0$. The relative nonzero class and the
spatial vanishing thus have explicit representatives
and comparison differentials.

\section{Spatial Hom and intersections of supports}
\label{cbc:ac:spatial-supports}

Full spatial Hom need not vanish for every
localization. Its answer depends on which support
intersections exist. Let $X=\mathbb A_k^d$,
$R=k[x_1,\ldots,x_d]$, and let
$H_i=\{\ell_i=0\}$ be a nonempty finite list of
affine hyperplanes defined over $k$. Use the full Weyl
algebra $D_X$ and the ordinary connection on
\[
 L=R[(\ell_1\cdots\ell_r)^{-1}].
\]

The \emph{intersection nerve} has a face
$S\subseteq\{1,\ldots,r\}$ when
$H_S=\bigcap_{i\in S}H_i$ is nonempty.
Put its augmentation in degree zero:
\begin{equation}\label{cbc:ac:nerve-complex}
 B^0=k,\qquad
 B^q=\bigoplus_{\substack{|S|=q\\H_S\ne\varnothing}}k
                       \quad(q\geq1).
\end{equation}
The differential is alternating incidence in
increasing index order. In particular,
$B^0\to B^1$ is the diagonal map.

\begin{theorem}\label{cbc:ac:spatial-nerve}
There is a natural equivalence
\[
 \operatorname{RHom}_{D_X}(L,R)\simeq B^\bullet
\]
that retains the support-inclusion maps of the
finite intersection diagram.
\end{theorem}
\begin{proof}
First let $W\subseteq X$ be a nonempty affine
linear subspace of codimension $c$. Choose
independent normal coordinates $u_1,\ldots,u_c$.
The local-cohomology complex $R\Gamma_W R$ is
the tensor of $[R\to R[u_i^{-1}]]$.
Thus it is $\delta_W[-c]$, where $\delta_W$
is the tensor of the normal principal-part
modules and the regular tangential module.

In one normal coordinate the free resolution of
the principal-part module is
$[D\xrightarrow{r_u}D]$.
Hom into $k[u]$ has cohomology in degree one,
and evaluation at zero identifies its cokernel
with $k$. Tangential polynomial de Rham
complexes contract to constants. Tensoring gives
\[
 \operatorname{RHom}_{D_X}(R\Gamma_W R,R)\simeq k.
\]

The canonical map to this complex from
$\operatorname{RHom}_{D_X}(R,R)\simeq k$
is a quasi-isomorphism. In one normal variable
the localization triangle proves this because
the reverse Hom from $k[u,u^{-1}]$ is
contracted by $(u\partial_u+1)^{-1}$.
Tensoring proves the assertion in all normal
directions.

For every inclusion of supports, these canonical
maps from the same complex
$\operatorname{RHom}_{D_X}(R,R)$ commute with
the induced trace map. They therefore give an
objectwise equivalence of diagrams, not only
separate identifications of their cohomology.
Every nonempty intersection contributes $k$,
each incidence coefficient is one, and an
empty intersection contributes zero.

Iterating the finite localization complexes
gives the Mayer--Vietoris complex for support
in $\bigcup_iH_i$. Applying derived Hom and
the preceding diagram equivalence gives the
ordinary cochains of the intersection nerve.
The triangle
$R\Gamma_{\cup_iH_i}R\to R\to L$ identifies
reverse Hom with the fiber of the augmentation
from $k$ to those cochains.
This is precisely \eqref{cbc:ac:nerve-complex}.
\end{proof}

When all hyperplanes meet, the nerve is a simplex.
Inserting a fixed vertex contracts its augmented
cochains. In particular, any nonempty arrangement
of equality hyperplanes has zero reverse Hom of
this localization into the regular target.

\begin{example}[An affine triple]\label{cbc:ac:affine-triple}
For $x=0$, $y=0$, and $x+y=1$, every pair
intersects and the triple does not. The complex is
\[
 k\xrightarrow{(1,1,1)}k^3
       \xrightarrow{(a,b,c)\mapsto(b-a,c-a,c-b)}k^3.
\]
It retracts onto $k[-2]$. With coordinates
$(w_{12},w_{13},w_{23})$ on the last term, take
\[
 \pi(w)=w_{12}-w_{13}+w_{23},\qquad
 i(\lambda)=(0,0,\lambda),
\]
and
$h_1(a,b,c)=a$, $h_2(w)=(0,w_{12},w_{13})$.
Substitution gives $dh+hd=1-i\pi$ and
$h^2=hi=\pi h=0$. Thus a full spatial reverse
class survives when the incidence diagram has
this missing intersection.
\end{example}

\begin{corollary}\label{cbc:ac:general-position}
Suppose $r>d\geq1$ hyperplanes are in affine
general position. This means that every subset
of at most $d$ intersects in the expected codimension
and every $(d+1)$-subset has empty intersection.
Then
\[
 \operatorname{RHom}_{D_X}(L,R)
                 \simeq k^{\binom{r-1}{d}}[-d].
\]
\end{corollary}
\begin{proof}
The nerve is the $(d-1)$-skeleton of the simplex.
For a degree-$q$ cochain $c$, set
\[
 (h_qc)_T=
 \begin{cases}
  c_{\{1\}\cup T},&1\notin T,\\
  0,&1\in T,
 \end{cases}
 \qquad |T|=q-1.
\]
Choose top representatives supported on the
$d$-subsets avoiding vertex one.
For $S=\{s_1<\cdots<s_d\}$ of this kind, put
\[
 (\pi_dc)_S=c_S-\sum_{j=1}^d(-1)^{j-1}
                       c_{\{1\}\cup(S\setminus\{s_j\})}.
\]
Alternating incidence gives
$dh+hd=1-i_d\pi_d$, with
$h^2=hi_d=\pi_dh=0$.
The number of these representatives is
$\binom{r-1}{d}$.
\end{proof}

\subsection{The global regular and localized pair}

For $R_n=k[z_1,\ldots,z_n]$ and
$\Delta_n=\prod_{i<j}(z_i-z_j)$, all equality
hyperplanes meet. Hence, for $n\geq2$,
\[
 \operatorname{RHom}_{D_n}
                (R_n[\Delta_n^{-1}],R_n)=0.
\]
At three points this vanishing determines the
marked pair consisting of the regular and fully
localized coefficient on the whole position space.

Put $R=R_3$ and $L=R[\Delta_3^{-1}]$. Define
\[
 E_3=\Lambda(\rho)\otimes_k(k\oplus ka\oplus kb),
 \qquad |\rho|=|a|=|b|=1,
\]
with zero differential and every product of two
elements of $ka\oplus kb$ equal to zero.
The tensor product uses the graded sign.

\begin{proposition}\label{cbc:ac:global-pair}
The derived endomorphism algebra of the marked
pair $(R,L)$ is formal. A thirteen-dimensional
model retaining both idempotents and the
localization map is
\[
 T(E_3)=
 \begin{pmatrix}k&0\\ E_3&E_3\end{pmatrix},
 \qquad d=0.
\]
The lower-left copy models
$\operatorname{RHom}(R,L)$, its element $1$
is localization, and composition is graded
matrix multiplication.
\end{proposition}
\begin{proof}
On the fully localized ring use
$c=z_3$, $u=z_1-z_3$, and
$t=(z_2-z_3)/(z_1-z_3)$. Then
\[
 L=k[c]\otimes k[u,u^{-1}]
          \otimes k[t,t^{-1},(t-1)^{-1}].
\]
There is a dg-algebra map from $E_3$ to
the algebraic de Rham complex of $L$ given by
\[
 \rho\longmapsto du/u,\qquad
 a\longmapsto dt/t,\qquad
 b\longmapsto dt/(t-1).
\]

It is a quasi-isomorphism. Polynomial integration
contracts the $c$ factor. In the $u$ factor,
only $1$ and $du/u$ survive. Partial fractions
in $t$ leave the two simple-pole forms.
Polynomial forms and higher poles have
primitives in the same ring, while residues
prove independence of the two remaining classes.
The two shape forms have zero product, so this
also verifies the multiplicative assertion.

Let $D_L$ be the localized differential-operator
algebra. Exact localization and its restriction
adjunction identify
$\operatorname{REnd}_{D_3}(L)$ with
$\operatorname{REnd}_{D_L}(L)$.
To compute its product, use the Spencer resolution
\[
 P=D_L\otimes_k\Lambda(e_1,e_2,e_3),\qquad
 |e_i|=-1,\qquad
 d_P=\sum_i r_{\partial_i}\iota_i
\]
in coordinates $(c,u,t)$, where $\iota_i$
contracts with $e_i$. Its order filtration
reduces exactness to the Koszul complex of
the derivative symbols.

There is a unital dg-algebra map
\[
 \Omega_L^\bullet\longrightarrow
                       \operatorname{End}_{D_L}(P),
 \qquad f\longmapsto r_f,\qquad dx_i\longmapsto-\iota_i.
\]
Indeed
$[d_P,r_f]=-\sum_i r_{\partial_i f}\iota_i$,
and the contractions anticommute and commute
with multiplication by functions.
Composition with augmentation to
$\operatorname{Hom}_{D_L}(P,L)$ is the de Rham
complex with these signs. The augmentation is
a quasi-isomorphism because $P$ is bounded
and projective. Thus the displayed algebra
map is a quasi-isomorphism.

The coordinate Spencer resolution of $R$
identifies $\operatorname{RHom}_{D_3}(R,L)$
with the same de Rham complex. Precomposition
with localization is the identity in these
models. Polynomial de Rham contraction gives
$\operatorname{REnd}_{D_3}(R)\simeq k$, and
the reverse corner vanishes by
Theorem~\ref{cbc:ac:spatial-nerve}.

To retain composition, choose a dg enhancement
of the pair and a closed localization map
$\lambda:R\to L$. Write $A=\operatorname{End}(L)$
there. The strict dg functor from $T(A)$
to this pair uses the scalar unit on
$\operatorname{End}(R)$, the identity on $A$,
and $a_0\mapsto a_0\lambda$ on the forward
corner. Its reverse corner is zero.
Every corner map is a quasi-isomorphism by
the calculations above. The multiplicative
zigzag $A\simeq\Omega_L^\bullet\leftarrow E_3$
therefore gives $T(A)\simeq T(E_3)$, preserving
the idempotents and localization map.
\end{proof}

On the chart $u\ne0$, the regular object becomes
$R[u^{-1}]$. The chart pair, viewed on the
whole space by restriction and direct image,
is consequently $(R[u^{-1}],L)$.
It has the nonzero reverse corner and ternary
obstruction of Theorem~\ref{cbc:ac:marked-nonformality}.
The whole-space pair in
Proposition~\ref{cbc:ac:global-pair} has regular
object $R$. These different objects account
for their different formality statements.
Neither statement by itself determines the
full exterior system with its primitive,
second exterior coefficient, diagonal
attachments, and vacuum maps.

\section{Higher braid transport}
\label{cbc:ac:braids}

The formal star algebra $A_r$ has no products between
its nonidentity arrows. It can nevertheless carry
transport with nonzero higher components. We construct
such transport on its bar, keeping the leaf generators
as part of the data.

Fix $r\geq2$ and put $S=\bigoplus_{i=0}^rke_i$.
The augmentation $A_r\to S$ kills every arrow.
Its reduced relative bar coalgebra is
\[
 BA_r=T_S^c((A_r/S)[1]),
\]
where $s$ lowers degree by one and $T^c$ is the
direct sum over tensor lengths, with deconcatenation
coproduct. Every suspended arrow has degree zero.
Thus $BA_r$ is concentrated in degree zero and
its bar differential is zero.

The full dual is the path-length completion
$\widehat P$ of the reverse quiver. Denote its
arrows by
\[
 x_i:0\longrightarrow i,\qquad
 y_i:i\longrightarrow0,
\]
where $x_i$ pairs with $s\beta_i$ and $y_i$
with $s\alpha_i$. Pairing reverses the order
of the tensors. There is no Koszul sign because
the suspended arrows have degree zero.

Products are composition with the rightmost map
applied first. The central loops and units are
\[
 z_i=y_ix_i,\qquad g_i=e_0+z_i,\qquad
 g_i^{-1}=\sum_{m\geq0}(-z_i)^m.
\]
Every finite path quotient is finite-dimensional.
Infinite expressions below mean their compatible
values in these quotients.

\begin{theorem}\label{cbc:ac:braid-substitution}
For $j=i+1$, exchange $e_i,e_j$, fix the other
idempotents, and define
\begin{align}
 \sigma_i(x_i)&=x_jg_i^{-1},&
 \sigma_i(y_i)&=g_i y_j,\notag\\
 \sigma_i(x_j)&=x_i,&
 \sigma_i(y_j)&=y_i.
                 \label{cbc:ac:braid-formulas}
\end{align}
Fix all other arrows. These substitutions are
continuous automorphisms of $\widehat P$.
They satisfy the braid relations and preserve
the ordered product $g_1\cdots g_r$.
Their inverses are
\[
 \begin{aligned}
  \sigma_i^{-1}(x_i)&=x_j,&
  \sigma_i^{-1}(y_i)&=y_j,\\
  \sigma_i^{-1}(x_j)&=x_i g_j,&
  \sigma_i^{-1}(y_j)&=g_j^{-1}y_i.
 \end{aligned}
\]
\end{theorem}
\begin{proof}
Every substituted arrow has positive length,
so the formulas define continuous endomorphisms.
Substitution verifies both inverse identities.
On the central units they give
\begin{equation}\label{cbc:ac:artin}
 \sigma_i(g_i)=g_ig_jg_i^{-1},\qquad
 \sigma_i(g_j)=g_i.
\end{equation}
Thus $g_ig_j$ and the full ordered product are
preserved.

For three consecutive leaves, write the central
units as $g,h,k$. Both sides of the adjacent
braid relation send the three arrow pairs to
\[
 (x_3h^{-1}g^{-1},ghy_3),\qquad
 (x_2g^{-1},gy_2),\qquad
 (x_1,y_1).
\]
They have the same permutation of vertices.
The relation therefore holds on all generators.
Distant substitutions act on disjoint pairs
and fix the central units used by the other
substitution, so they commute.
These identities hold in words in invertible
units before their series expansions.
\end{proof}

The central formulas \eqref{cbc:ac:artin} are
the usual Artin substitutions. Formula
~\eqref{cbc:ac:braid-formulas} lifts them to
the arrows themselves. This lift retains a
chosen generator at each leaf, which is the
framing of the path action.

Transpose $\sigma_i$ against the finite path
spaces. This gives a coaugmented coalgebra
automorphism $F_i$ of $BA_r$, with its stated
permutation of objects. Its corestrictions
are specified in every arity by
\begin{equation}\label{cbc:ac:braid-transpose}
 \langle q,F_{i,n}(sa_1|\cdots|sa_n)\rangle
      =\langle\sigma_i(q),sa_1|\cdots|sa_n\rangle
\end{equation}
for each dual arrow $q$.

\begin{proposition}\label{cbc:ac:braid-higher}
The components \eqref{cbc:ac:braid-transpose},
extended by strict unitality, define augmented
$A_\infty$-automorphisms of $A_r$. They satisfy
the braid relations. Every even component
vanishes, while the expansion
\[
 \sigma_i(x_i)=\sum_{m\geq0}(-1)^m x_jz_i^m,\qquad
 \sigma_i(y_i)=y_j+z_i y_j
\]
gives nonzero components in every odd arity
in the first family.
\end{proposition}
\begin{proof}
Transposing multiplication gives the coalgebra
identity. The zero bar differential gives the
chain identity. On an input of fixed length,
only finitely many path coefficients can
contribute, since substitution never decreases
path length. Thus the transpose preserves
the direct-sum bar.

Inverses transpose to inverses. Transposition
reverses composition, but the adjacent braid
words are palindromes and distant commutation
is symmetric. The braid relations therefore
remain valid. Every path parallel to an
arrow has odd length. The displayed series
then determines the parity and nonvanishing
of the components. Adjoining the identity
operations gives strict unitality.
\end{proof}

For the pure braid $\Phi_i=\sigma_i^2$, the
vertices and linear terms are fixed. In
central loop length,
\begin{equation}\label{cbc:ac:pure-loops}
 \begin{aligned}
  \Phi_i(z_i)&=z_i-[z_i,z_j]+O(z^3),\\
  \Phi_i(z_j)&=z_j+[z_i,z_j]+O(z^3),\\
  \Phi_i(z_k)&=z_k\quad(k\notin\{i,j\}).
 \end{aligned}
\end{equation}
Here $O(z^3)$ means terms of loop length at
least three. On arrows the first corrections
include
\[
 \Phi_i(x_i)=x_i-x_i z_j+O(\text{length }5),\qquad
 \Phi_i(y_i)=y_i+z_j y_i+O(\text{length }5).
\]
Their transposes have identity linear term
and nonzero ternary terms.

\begin{proposition}\label{cbc:ac:pure-nontrivial}
For $r\geq3$, $\Phi_i$ is not inner on the
central completed free algebra. For $r=2$,
the full path action is conjugation by
\[
 u=g_1g_2+(e_1+x_1y_1)+(e_2+x_2y_2).
\]
In either case, its transpose is not
homotopic to the identity through augmented,
strictly unital bar maps with the objects
fixed.
\end{proposition}
\begin{proof}
An inner automorphism of the central corner
is conjugation by a unit with nonzero scalar
constant term. Its degree-two correction is
the commutator with a linear term
$h=\sum_k a_k z_k$.
Formula~\eqref{cbc:ac:pure-loops} forces
$a_i=a_j=1$ and every other $a_k=0$.
For a third generator, however,
$[h,z_k]=[z_i+z_j,z_k]\ne0$, while
$\Phi_i$ fixes $z_k$. This proves the first
assertion.

For two leaves, set $u_0=g_1g_2$ and
$u_j=e_j+x_jy_j$. These are units in their
corners and satisfy
$u_jx_j=x_jg_j$, $y_ju_j=g_jy_j$.
Conjugation by their sum sends
\[
 \begin{aligned}
 (x_1,y_1)&\longmapsto
   (x_1g_1g_2^{-1}g_1^{-1},\,g_1g_2g_1^{-1}y_1),\\
 (x_2,y_2)&\longmapsto(x_2g_1^{-1},\,g_1y_2).
 \end{aligned}
\]
These are the twofold substitutions of
\eqref{cbc:ac:braid-formulas}.

The reduced bars are concentrated in degree
zero. A degree-minus-one bar homotopy must
therefore vanish. Equivalently, a chain map
to the bar tensored with polynomial interval
forms has no degree-minus-one coefficient
for a $dt$ term. Its degree-zero component
has derivative zero and is constant.
The nonzero ternary term rules out such a
homotopy to the identity.
\end{proof}

The two-leaf action is inner, but that inner
conjugation is not an augmented bar homotopy.
The equivalence relation is part of the
transport statement.

\section{Logarithmic families and their curvature}
\label{cbc:ac:logarithmic}

Let $\Phi$ be a product of the pure-braid
substitutions. The operator $\Phi-1$ raises
arrow length by at least two. Its logarithm
therefore exists on every finite path quotient:
\[
 \mathcal L=\log\Phi
        =\sum_{n\geq1}\frac{(-1)^{n+1}}n(\Phi-1)^n.
\]
The resulting operators are compatible under
quotient maps.

\begin{lemma}\label{cbc:ac:log-derivation}
The operator $\mathcal L$ is a continuous
derivation. It raises arrow length by at
least two and annihilates every vertex
idempotent and $g_1\cdots g_r$.
\end{lemma}
\begin{proof}
At any finite path precision, the expression
\[
 \Phi^v=\sum_{n\geq0}\binom vn(\Phi-1)^n
\]
is a polynomial in $v$. For every nonnegative
integer $v$ it is multiplicative. In
characteristic zero this is a polynomial
identity, so it holds for arbitrary $v$.
Differentiating at zero proves the Leibniz
rule for $\log\Phi$.
The length assertion follows termwise.
Every element fixed by $\Phi$ is killed
by each term of its logarithm, which proves
the remaining assertions.
\end{proof}

Take $k=\mathbb C$ and write
$D=\mathcal L/(2\pi i)$. On a punctured
disc with coordinate $q$, define
\begin{equation}\label{cbc:ac:log-connection}
 \nabla=d-D\,\frac{dq}{q}.
\end{equation}
At a finite path quotient $D$ is nilpotent.
Horizontal transport is
$\exp(D\log q)$, a finite polynomial in
$\log q$, with counterclockwise monodromy
$\Phi$.

Each quotient has the corresponding
holomorphic lattice with a logarithmic
connection at $q=0$. These lattices form
a compatible inverse system. The same
formula defines a formal logarithmic
connection over $\mathbb C[[q]]$.

\begin{theorem}\label{cbc:ac:log-family}
Connection~\eqref{cbc:ac:log-connection} is
flat and multiplicative. The paired bar
coalgebra has connection
\[
 \nabla^\vee=d+D^\vee\,\frac{dq}{q}
\]
and monodromy $\Phi^{-\vee}$.
It defines a strictly unital augmented
$A_\infty$-algebra over logarithmic de Rham
coefficients. Its new reduced bar components
are
\[
 b_n^{\mathrm{new}}=\frac{dq}{q}(D^\vee)_n,
                         \qquad n\geq3.
\]
Only odd arities occur. Every point
evaluation recovers $A_r$.
\end{theorem}
\begin{proof}
The derivation property of $D$ proves
multiplicativity. The form $dq/q$ is
closed and has square zero, so the
connection is flat. The plus sign on
the dual makes evaluation horizontal
and gives inverse-transpose monodromy.

The transpose of a derivation is a
coderivation. Here it lowers length by
at least two, so it acts finitely on
each bar input and has no unary
component. Path parity excludes even
arities. On the bar tensored with
logarithmic de Rham coefficients,
$d+(dq/q)D^\vee$ is a degree-one
coderivation of square zero. This is
the $A_\infty$ bar equation.

After desuspension, the arity-$n$
operation has degree $2-n$.
The original unit and binary product
remain, and all new operations vanish
when an input is a unit.
Point evaluation kills $dq/q$.
\end{proof}

The pointwise algebra is constant while
its higher transport is nontrivial.
The boundary product is horizontal by
Lemma~\ref{cbc:ac:log-derivation}.
These statements concern the completed
path algebra and its paired bar.

Two collision parameters impose an
additional equation. For three leaves
put $\mathcal L_{12}=\log(\sigma_1^2)$
and $\mathcal L_{23}=\log(\sigma_2^2)$.

\begin{proposition}\label{cbc:ac:log-curvature}
The two derivations have nonzero commutator:
\[
 [\mathcal L_{12},\mathcal L_{23}](x_1)
    =x_1(z_3z_2-z_2z_3)+O(\text{length }7).
\]
Consequently, on a product of two punctured
discs the connection
\[
 d-\frac{\mathcal L_{12}}{2\pi i}\,d\log q_1
   -\frac{\mathcal L_{23}}{2\pi i}\,d\log q_2
\]
is not flat.
\end{proposition}
\begin{proof}
The first nonzero terms of a logarithm
agree with those of its pure map minus
the identity. Thus
\[
 \mathcal L_{12}(x_1)=-x_1z_2+O(\text{length }5),
 \qquad \mathcal L_{23}(x_1)=0,\qquad
 \mathcal L_{23}(z_2)=[z_3,z_2]+O(z^3).
\]
The derivation rule gives the stated
commutator. Its leading paths are
distinct, so their difference is nonzero.
The curvature is
\[
 (2\pi i)^{-2}[\mathcal L_{12},\mathcal L_{23}]
                   \,d\log q_1\wedge d\log q_2.
\]
\end{proof}

For commuting pure monodromies the logarithms
commute, and the sum of their one-variable
connections is flat. The displayed adjacent
pure monodromies instead require further
terms satisfying the curvature equation
if they are to belong to a simultaneous
family. One-variable transport alone
does not construct those terms.

The analytic line gauge and the framed
path action were constructed on specified
coefficient spaces. To identify either
with native collision transport requires
maps from the supported collision and
endpoint modules. Such maps must preserve
the normal derivatives, the regular
attachments, both vacuum insertions, and
nested compositions on a common continuous
domain. The separated estimates supply
that domain for the line calculation.
The finite marked models specify which
objects and higher operations a collision
comparison must retain. An identification
with the full ordered collision atlas
requires this further construction.
